\chapter{Rizq as Conserved Quantity}
\label{ch:rizq-conservation}

\epigraph{\itshape No soul will leave this world until its provision
reaches it in full.}{--- \hadith{Sahih Muslim} 2664}

\noindent
This chapter is the first of Part III. The eight chapters of Part III
take the canonical Islamic metaphysical claims and treat each
rigorously, using the methodological apparatus of Part I and the
analytical apparatus of Part II. Each chapter follows the same
structure: state the claim, present the skeptical objection,
formalize the claim using the modern apparatus, derive its content,
state the empirical predictions, acknowledge what remains open, and
identify the level of treatment achieved (\xref{ch:levels-of-rigor}).

The present chapter treats the claim that \arb{rizq} is conserved
over a lifetime: that no soul dies before completing its allotted
provision. The treatment is Level 2 in our taxonomy, with elements
moving toward Level 3 where the apparatus permits. The chapter
establishes the formal framework; subsequent chapters will use the
same framework with modifications.

The chapter proceeds in nine sections.
\xref{sec:rizq-claim} states the scriptural claim precisely and
catalogs its loci.
\xref{sec:rizq-objection} presents the skeptical objection.
\xref{sec:rizq-formalization} formalizes the claim as a conservation
law in the sense of Noether's theorem.
\xref{sec:rizq-stochastic} develops the stochastic-process formulation.
\xref{sec:rizq-bohmian} treats the Bohmian interpretation in which
the lifetime trajectory is fully determined.
\xref{sec:rizq-empirical} discusses what empirical evidence
would support or refute the claim.
\xref{sec:rizq-akbarian} integrates the Akbarian framework of
\xref{ch:akbarian}.
\xref{sec:rizq-implications} draws the practical and theoretical
implications.
\xref{sec:rizq-open} states what remains open.

\section{The scriptural claim}
\label{sec:rizq-claim}

\subsection{The Quranic basis}

We collected the Quranic loci in \xref{sec:quran}; we restate the
most central here. The locus classicus is \Quran{11:6}:

\begin{quote}
\textit{And there is no creature on earth but that upon Allah is its
provision (\arb{rizq}), and He knows its place of dwelling and
place of storage. All is in a clear register.}
\end{quote}

The verse asserts three propositions: (1) every living creature has
provision that Allah has assumed responsibility for; (2) the place
and manner of provision are known to Allah; (3) all of this is
recorded in a clear register. The classical exegesis (al-Tabari,
al-Razi, al-Qurtubi) is unanimous that ``upon Allah''
(\arb{`ala Allah}) carries the full force of divine undertaking.

The verse \Quran{29:60} extends the claim beyond humans:

\begin{quote}
\textit{And how many a creature carries not its own provision!
Allah provides for it and for you.}
\end{quote}

The decoupling of provision from the creature's provisioning capacity
is explicit. The wild animal that cannot save or plan still receives
its provision; the same is true, the verse implies, for human beings.

The temporal structure is supplied by the hadith we cite as the
chapter's epigraph and develop now.

\subsection{The hadith on the lifetime integral}

\hadith{Sahih Muslim} 2664 states:

\begin{quote}
\textit{No one of you will leave this world until his provision
(\arb{rizq}) reaches him in full and his appointed term
(\arb{ajal}) has been completed.}
\end{quote}

The hadith makes two claims simultaneously, and the simultaneity is
what we must formalize. First, the lifetime integral of provision is
fixed; it ``reaches in full'' before death. Second, the time of death
is similarly fixed. Together they imply that no individual dies
having received less than their decreed provision.

A second key report from \hadith{Sahih al-Bukhari} 3208 establishes
the pre-natal recording of \arb{rizq}:

\begin{quote}
\textit{...Then Allah sends an angel and orders him to write four
things: his provision (\arb{rizq}), his lifespan (\arb{ajal}),
his deeds, and whether he will be among the wretched or the
felicitous.}
\end{quote}

The hadith establishes that the decree of provision is not merely
chronologically prior to the soul's life; it is part of the
ontological structure within which the soul is created. The
provision is recorded before the soul exists.

\subsection{The composite claim}

We can now state the composite claim that the chapter will formalize:

\begin{conceptbox}[The rizq-conservation claim]
\label{conc:rizq-claim}
For each created soul $i$, there exist:
\begin{itemize}[leftmargin=1.5em]
    \item A fixed total provision $K_i \in \mathbb{R}_+$,
    pre-determined.
    \item A fixed lifespan $T_i \in \mathbb{R}_+$, pre-determined.
\end{itemize}
The instantaneous provision rate $R_i(t)$ for soul $i$ at time $t$
satisfies the conservation condition
\[
    \int_0^{T_i} R_i(t)\, dt = K_i.
\]
The lifetime integral is invariant; the temporal distribution
$R_i(t)$ may vary, but the integral is fixed.
\end{conceptbox}

This is the precise claim. The rest of the chapter examines what it
means, what apparatus formalizes it, what implications follow, and
what evidence would bear on it.

\section{The skeptical objection}
\label{sec:rizq-objection}

The skeptical reader will object on several grounds. We treat the
strongest formulations.

\subsection{The free-will objection}

If $K_i$ and $T_i$ are pre-determined, then the soul's striving has
no effect on the integral. Why work? Why save? Why plan? The
classical \arb{kalam} debate on \arb{qadar} (\xref{sec:kalam}) is
the historical articulation of this objection. We treat it
fully in \xref{ch:qadar} (Chapter 16). For the present chapter, we
note that the objection applies equally to any deterministic
framework (Newtonian physics included) and is dissolved by the
Bohmian-style structure we develop in \xref{sec:rizq-bohmian}.

\subsection{The empirical objection}

The objection: where is the empirical evidence? Surely we have
abundant data on individual provision over lifetimes; if the
integral is invariant, this should be observable. The data, as far
as we know it, does not show such an invariance. People starve.
People die in want. The rich and the poor experience radically
different lifetime integrals. How can the claim be true?

This objection rests on a misreading of what is being claimed. The
$K_i$ in the conservation claim is not a constant across souls;
each soul has its own $K_i$, decreed in advance. The claim is not
that all souls receive the same provision; it is that each soul
receives the provision allotted to it. The differences across
individuals are not evidence against the conservation claim; they are
expressions of differential decrees.

The empirical question is therefore harder than it first appears.
What we would need to establish, to confirm the claim, is that for
each individual the lifetime integral is invariant under
counterfactual variations of the individual's life history.
Different choices, different circumstances, different efforts ---
do they all produce the same integral? This is what the claim
asserts, and it is much harder to test than the cross-individual
question.

We discuss the empirical strategy in \xref{sec:rizq-empirical}.

\subsection{The metaphysical objection}

The objection: how can the integral of a temporal flow be
pre-determined before the flow occurs? The provision arrives
day by day; what does it mean for the total to be fixed in advance?

This objection is the deepest. It can be answered in several ways.
The Akbarian framework of \xref{ch:akbarian} provides one answer:
the soul itself, with its specific structural capacity to receive
the manifestation of \arb{ar-Razzaq}, is what has been determined;
the integral $K_i$ is the structural property of this capacity. The
Bohmian framework of \xref{sec:rizq-bohmian} provides another:
the trajectory of the soul through the configuration space of
possible life histories is fully determined by the initial
conditions, and the integral is a property of that trajectory. We
develop both.

\section{Formalization as a conservation law}
\label{sec:rizq-formalization}

\subsection{The Noether structure}

The conservation claim has the formal structure of a conservation
law. In physics, conservation laws are not accidents; by Noether's
theorem (1918), every continuous symmetry of the Lagrangian
corresponds to a conserved quantity. Conservation of energy follows
from time-translation symmetry; conservation of momentum follows
from spatial-translation symmetry; conservation of charge follows
from gauge symmetry.

If we treat the claim seriously as a Level 2 formal analogy
(\xref{ch:levels-of-rigor}), we should ask: what is the symmetry
that produces the conservation of $K_i$?

\begin{conceptbox}[Candidate symmetry: counterfactual life-path invariance]
\label{conc:counterfactual-invariance}
Let $\Gamma_i$ denote the set of all possible life-paths of soul
$i$ --- all the trajectories the soul could have taken, given its
initial conditions and the decree governing its existence. For each
$\gamma \in \Gamma_i$, let $R_i^\gamma(t)$ denote the provision rate
along path $\gamma$, and let
\[
    K_i^\gamma = \int_0^{T_i} R_i^\gamma(t)\, dt.
\]
The conservation claim is: $K_i^\gamma = K_i$ for all $\gamma \in
\Gamma_i$. The lifetime integral is invariant under the
transformation ``take a different counterfactual life path.''
\end{conceptbox}

This is a strong claim with the formal structure of a Noether
invariance. If we further suppose that $\Gamma_i$ is a continuous
family parameterized by some variable (the soul's degree of
striving, perhaps), then the invariance corresponds to a
continuous symmetry, and Noether's theorem implies a conserved
quantity. The conserved quantity is $K_i$.

\subsection{What this formalization gives us and does not give us}

The Noether-type formalization gives us a precise structural claim.
It tells us that the invariance of $K_i$ has the same structural
character as the invariance of energy in physics. The conservation
is not arbitrary; it expresses a symmetry of the underlying
structure.

The formalization does not give us a physical mechanism. We are not
claiming that the Lagrangian of the soul is well-defined or that
Noether's theorem applies in its physical sense. We are claiming
that the claim has the structural form of a conservation law and
that this form admits Noether-type formal treatment. This is
exactly Level 2 in our taxonomy: a formal analogy that imports
structural insight without commitment to the underlying mechanism.

\subsection{What the symmetry is symmetry of}

A subtle point. In physics, the symmetry that produces a conservation
law is a symmetry of the action functional or the Hamiltonian, both
of which are mathematical objects. In the present case, what is the
symmetry of?

The natural answer, within the Akbarian framework, is that the
symmetry is a symmetry of the soul's structural capacity. The soul,
as a locus of divine self-disclosure (\xref{sec:divine-names}), has a
specific structural character that determines what manifestations it
can receive. The integral $K_i$ is the magnitude of the soul's
capacity for receiving the manifestation of \arb{ar-Razzaq}.
Different counterfactual life paths --- different temporal
distributions of the manifestation --- do not change the soul's
structural capacity, and therefore do not change the integral.

The symmetry, in Akbarian terms, is the invariance of the soul's
capacity under variation of the temporal manifestation of that
capacity. The capacity is what is fixed; the manifestation is what
varies. The conservation is the natural consequence.

\section{The stochastic-process formulation}
\label{sec:rizq-stochastic}

\subsection{Provision as a stochastic process}

Independent of the Noether-type formalization, we can formalize the
conservation claim as a property of stochastic processes. Let
$R_i(t)$ be a stochastic process representing soul $i$'s provision
rate at time $t$. The conservation claim is:

\begin{conceptbox}[Stochastic conservation]
\label{conc:stochastic-conservation}
The lifetime integral $\int_0^{T_i} R_i(t)\, dt$ is, with probability
1 in the appropriate measure, equal to $K_i$. The process $R_i(t)$
varies stochastically, but the integral is deterministic.
\end{conceptbox}

This is the time-integral analogue of the deterministic-but-real
structure of Bohmian mechanics that we treat in \xref{sec:rizq-bohmian}.
The process is random in its instantaneous values; its integral is
not.

\subsection{Mathematical implications}

This formulation has specific mathematical implications. Standard
stochastic processes (Brownian motion, Poisson processes) do not
generally have deterministic integrals. The variance of the
integral typically grows with the time horizon. The
rizq-conservation claim asserts that the variance of the integral
is zero.

There exist mathematical objects with this property. They are
sometimes called \emph{constrained stochastic processes}: stochastic
processes whose paths are restricted to satisfy a deterministic
constraint. The standard mathematical machinery for such processes
is well-developed in the theory of stochastic differential equations
with constraints.

The simplest model: a process $R_i(t)$ that is otherwise random but
whose accumulated integral is monitored, with adjustments to keep
the integral on track for the predetermined value $K_i$. Such a
process has random-looking paths (the $R_i(t)$ values vary) but
deterministic integral (the $\int_0^{T_i} R_i(t)\, dt$ is fixed).

The technical apparatus is the apparatus of bridging stochastic
processes, particularly the Brownian bridge. A Brownian bridge is a
Brownian motion conditioned to take a specific value at a specific
endpoint; its paths are random but its endpoint is fixed. The
\arb{rizq} process can be modeled as an analogous bridge: a process
whose instantaneous values are random but whose terminal integral
is fixed.

\subsection{Implications for striving}

This formulation answers the free-will objection of
\xref{sec:rizq-objection}. The fact that $K_i$ is fixed does not
mean that $R_i(t)$ is fixed; the temporal distribution of provision
varies. Striving is real and effective; it shapes the
distribution of provision across the lifetime. What striving cannot
do is change the integral.

This is consistent with the classical hadith on \arb{tawakkul}:

\begin{quote}
\textit{If you all relied on Allah with the proper trust, He would
provide for you as He provides for the birds: they leave (their
nests) hungry in the morning and return filled in the evening.}
(\hadith{Tirmidhi} 2517)
\end{quote}

The hadith does not say humans should not work or plan. It asserts
that the integral is unaffected by the worry; what worry can do is
distort the temporal distribution. The believer who works without
worry has the same integral as the believer who works with worry,
but the work is more effective and the experience is more peaceful.

\section{The Bohmian formulation}
\label{sec:rizq-bohmian}

\subsection{The Bohmian structure applied to lifetime trajectories}

Bohmian mechanics (\xref{sec:interpretations}) gives us a formal
structure in which trajectories are fully determined yet motion is
real. We apply this structure to the soul's life trajectory.

\begin{conceptbox}[Bohmian formulation of rizq conservation]
\label{conc:bohmian-rizq}
The configuration space for soul $i$ is the space of all possible
life trajectories. The wave function $\Psi_i$ on this space is
determined by the divine decree at the moment of the soul's
creation. The actual life of the soul is a single trajectory through
the configuration space, guided by $\Psi_i$ according to a
guidance equation analogous to the Bohmian guidance equation
(\xref{prin:bohmian}).

The integral $K_i = \int_0^{T_i} R_i(t)\, dt$ is a functional of
the trajectory. Because the trajectory is fully determined by
$\Psi_i$ and the soul's initial position in configuration space,
the integral is fully determined as well. Yet the motion along the
trajectory is real; the soul actually traverses the trajectory.
\end{conceptbox}

This is a formal structural claim. We are not asserting that the
soul is a Bohmian particle. We are asserting that the structure
``determined trajectory + real motion + fixed integral'' is
mathematically coherent and physically instantiated in Bohmian
mechanics. Whatever the actual mechanism of the soul's life, it
can have this structure without contradiction.

\subsection{The dissolution of the determinism objection}

The free-will objection of \xref{sec:rizq-objection} is now
dissolved. In the Bohmian framework, the trajectory is fully
determined yet the motion is real. Striving is not illusory; it is
the soul's actual motion through the configuration space. The fact
that the motion is determined does not make it less real.

The classical \arb{kalam} debate on \arb{qadar}, treated in detail
in \xref{ch:qadar} (Chapter 16), can be read as the attempt to
articulate this structure in pre-modern vocabulary. The Mu`tazila,
Ash`ariyya, and Maturidiyya each tried to formulate a position that
preserved both divine determination and human responsibility. Each
formulation has its difficulties; all are addressing the same
underlying structure that Bohmian mechanics articulates cleanly.

The claim of the present chapter is therefore not just about
\arb{rizq}; it is about a class of phenomena that have the
structure ``determined integral + real path + variable
distribution.'' The structure is not paradoxical; it is exhibited
by physical systems that we now understand mathematically. The
metaphysical claim of the Islamic tradition is consistent with the
mathematical structure.

\subsection{Why this is Level 2 and not Level 3}

The treatment of this chapter is Level 2 in our taxonomy
(\xref{ch:levels-of-rigor}). We have given the conservation claim a
formal structure (Noether-type, stochastic-process, Bohmian) but we
have not identified the specific physical mechanism that produces
the conservation. The mechanism, in the Akbarian framework, is the
soul's structural capacity to receive divine self-disclosure
(\xref{sec:rizq-akbarian}); but ``the soul's structural capacity''
is itself a metaphysical claim, not a physical mechanism in the
modern scientific sense.

To reach Level 3 we would need to identify the physical
correlates of the soul's structural capacity --- some neural,
behavioral, or social variables whose variation across individuals
predicts variation in $K_i$. We do not currently have these. The
operational research program is one of the directions in which the
Hidden Architecture project must extend; we discuss it in
\xref{sec:rizq-empirical}.

\section{The empirical challenge}
\label{sec:rizq-empirical}

\subsection{What would constitute evidence}

The conservation claim is, in principle, empirical: $K_i$ is a real
quantity for each soul, and we can in principle measure
$\int_0^{T_i} R_i(t)\, dt$ for individuals whose lives have
concluded. The challenge is what would count as a test of the
claim.

The naive test --- compare $K_i$ across individuals --- does not
test the claim. Different individuals have different decreed
$K_i$; the cross-individual variation is consistent with the claim,
not evidence against it.

The test we would want is counterfactual: for a given individual,
would they have had the same $K_i$ if their life had unfolded
differently? This is impossible to test directly (we cannot rerun a
life), but it can be approached statistically through carefully
designed natural experiments.

\subsection{Natural experiments}

Several types of natural experiment are relevant.

First, twins reared apart. Identical twins share genetic endowment
but, in the case of separation, experience different environments
and different life histories. The conservation claim predicts that
their lifetime $K_i$ values should be more similar than would be
expected from random allocation among their shared environments.
This is a testable prediction.

Second, lottery winners. Individuals who win large lotteries
experience a sudden, exogenous increase in $R_i(t)$ at a particular
time. The conservation claim predicts that subsequent $R_i(t)$
should adjust downward to keep the integral on track. The
empirical literature on lottery winners (Imbens et al.\ 2001 and
subsequent work) shows that lottery winners typically reduce their
labor supply and experience a return to income trajectories not
dramatically different from non-winners; this is suggestive but
not definitive.

Third, sudden adverse events. Individuals who experience sudden
losses (asset seizures, business failures) might be expected to
have subsequent $R_i(t)$ adjustments upward to keep the integral on
track. The data on this is harder to interpret because the events
themselves are often correlated with personal characteristics that
affect subsequent income.

Fourth, longitudinal cohort studies. Long-running cohort studies
(the Framingham study, the British cohort studies, etc.) provide
detailed individual income data over decades. Statistical analysis
of the variance of lifetime income, conditional on early-life
predictors, can shed light on the conservation claim.

\subsection{Methodological challenges}

Several methodological challenges must be acknowledged.

First, defining $R_i(t)$ operationally. Is it monetary income? Is it
real consumption? Does it include non-monetary provision (housing
provided by family, food cultivated, etc.)? The choice of
operationalization affects the empirical results. The Akbarian
framework would suggest that $R_i(t)$ is the full manifestation of
\arb{ar-Razzaq} in the soul's life, which is broader than monetary
income. But broader operationalizations are harder to measure.

Second, the lifetime $T_i$ is not known until death. Empirical
studies of lifetime integrals are necessarily retrospective and
conducted on completed lives.

Third, the comparison group problem. To test the conservation claim
statistically, we need to compare actual lifetime integrals with
counterfactual ones. The counterfactual is unobservable.

Despite these challenges, the empirical program is not impossible.
The challenges are common to many areas of social science research
(the difficulty of counterfactual reasoning, operationalization
choices, comparison group construction) and methodologies have been
developed for addressing them. The conservation claim is empirically
testable in principle; the empirical program has not yet been
seriously undertaken; it is one of the contributions the Hidden
Architecture project can stimulate.

\section{The Akbarian integration}
\label{sec:rizq-akbarian}

\subsection{Rizq as manifestation of ar-Razzaq}

The Akbarian framework of \xref{ch:akbarian} provides the
metaphysical substrate for the conservation claim. \arb{Rizq} is
not an independent substance that gets distributed; it is the
manifestation of the divine name \arb{ar-Razzaq} (the Provider) in
the locus of the soul. The integral $K_i$ is the structural
capacity of the soul's locus for receiving this manifestation.

This metaphysical substrate explains several features of the formal
treatment.

\textbf{Why the integral is fixed.} The integral is fixed because
the soul's structural capacity is fixed. The capacity is not a
property the soul earns; it is what the soul is, ontologically.
The soul is created with a specific capacity for receiving the
manifestation of \arb{ar-Razzaq}, and the integral is just the
total magnitude of this capacity actualized over the lifetime.

\textbf{Why the temporal distribution varies.} The temporal
distribution can vary because the actualization of the capacity is
a temporal process subject to the soul's choices, the surrounding
circumstances, and the moment of the divine disclosure. The
capacity is fixed; the temporal pattern of its actualization is
not.

\textbf{Why striving matters but does not change the integral.}
Striving affects the temporal pattern of actualization. The soul
that strives is more receptive to the manifestation, more able to
recognize it when it arrives, more able to use it for the soul's
intended purposes. But the magnitude of the manifestation is fixed
by the capacity, not by the striving.

\textbf{Why the integral varies across souls.} Different souls have
different capacities. Different capacities correspond to different
configurations of divine names in the soul (\xref{sec:divine-names});
the integral of \arb{ar-Razzaq} manifestation in a soul whose
configuration emphasizes \arb{al-Ghani} (the Self-Sufficient) is
different from the integral in a soul whose configuration
emphasizes \arb{al-Wahhab} (the Constant Giver).

\subsection{The dynamics of tajalli}

The Akbarian doctrine of \arb{tajalli} (continuous self-disclosure,
\xref{sec:tajalli-mechanism}) gives the dynamic content of the
conservation claim. The provision $R_i(t)$ is the instantaneous
manifestation of \arb{ar-Razzaq} in the soul at time $t$. The
manifestation is continuous, not once-and-for-all; the soul receives
fresh provision moment by moment, throughout the lifetime.

The integral over the lifetime is the total magnitude of the
manifestation. The total is fixed by the soul's capacity. The
moment-by-moment distribution is the temporal unfolding of the
manifestation, subject to the soul's actual life history.

This is a richer picture than the bare conservation claim. It
embeds the claim in a metaphysical framework in which the
conservation is not a strange constraint imposed from outside but
a natural consequence of the structural relationship between the
soul and the divine self-disclosure.

\subsection{Comparison with classical theological treatments}

The classical theological treatment of \arb{rizq} in the
mainstream tradition emphasized the divine guarantee of provision
and the futility of excessive worry, but did not articulate the
conservation structure formally. Texts like al-Ghazali's
\textit{Ihya'} and later devotional literature treat \arb{rizq} as
a divine assurance to be trusted; the formal mathematical structure
we have given here was not made explicit.

This is one of the contributions of the present project. The
classical treatment had the substance right; what was missing was
the formal framework that allows the substance to be expressed in
modern theoretical vocabulary. The conservation structure, the
stochastic-bridge formalism, the Bohmian analogue --- these
articulate, in modern terms, what the classical tradition was
asserting in its own vocabulary.

\section{Implications}
\label{sec:rizq-implications}

\subsection{Practical implications}

The conservation claim has direct practical implications for the
believer.

\textbf{Anxiety reduction.} If the integral is fixed by the
soul's capacity, then anxiety about the future is misplaced. The
provision will arrive; the only question is when and how. The
reduction of anxiety is itself a behavioral consequence with
empirical correlates; we treat it in \xref{ch:gratitude}.

\textbf{Effort allocation.} If striving affects the temporal
distribution but not the integral, then the optimal effort is the
effort that makes the temporal distribution match the soul's
purposes, not the effort that maximizes the integral. This is
substantively different from the standard maximization framework
and has implications for retirement planning, savings behavior,
and intergenerational wealth transfer.

\textbf{Charity behavior.} If the integral is fixed, then charity
does not reduce the giver's lifetime provision; it shifts the
temporal distribution. This is consistent with the Quranic
teaching that charity does not deplete wealth (\Quran{2:261},
\Quran{34:39}, and the hadith corpus).

\subsection{Theoretical implications}

The conservation claim has implications for several other parts of
the Hidden Architecture project.

\textbf{For Chapter 11 (network of provision).} The claim that
children bring their own provision is, in part, the claim that
each child has its own conservation integral $K_i$. The household-network
analysis of \xref{ch:network-provision} builds on this individual
claim.

\textbf{For Chapter 14 (gratitude).} The behavioral mediation chain
through which gratitude affects long-run wealth includes the
disposition that recognizes provision as it arrives. The
recognition is what allows the soul to receive its $K_i$ fully
within the lifetime $T_i$.

\textbf{For Chapter 15 (barakah).} The barakah multiplier on
material wealth is, on the Akbarian reading, the intensity with
which the manifestation of \arb{ar-Razzaq} is occurring. Barakah
does not change the integral; it changes the experiential
character of the provision.

\textbf{For Chapter 16 (qadar).} The conservation claim is a
specific instance of the general qadar structure. The Bohmian and
QBist treatments of qadar in Chapter 16 apply directly here.

\subsection{Implications for economic policy}

If the conservation claim is approximately correct, several
mainstream economic policies are at minimum questionable.
Aggressive growth maximization, which assumes that economic policy
can change individual lifetime provision integrals, may be
mistaking distributional shifts for integral increases.
Anti-poverty policies that focus on the lifetime integral may be
less effective than those that focus on smoothing the temporal
distribution.

This is a direction that subsequent volumes of the Islamic
Economics series will develop. The present chapter establishes the
foundational claim; the policy implications are downstream.

\section{What remains open}
\label{sec:rizq-open}

The conservation claim, as we have formalized it, is at Level 2. To
reach Level 3 we would need:

\begin{enumerate}[leftmargin=2em]
    \item An operational specification of $R_i(t)$ that makes it
    measurable across individuals and time.
    \item Empirical data, from natural experiments or longitudinal
    studies, that supports or refutes the conservation claim.
    \item A specification of the physical correlates of the soul's
    ``structural capacity'' that bears on the integral.
\end{enumerate}

None of these is currently available. The empirical program for
the rizq-conservation claim is an open research direction.

We should also acknowledge limits of the formal treatment. The
Bohmian formulation is structural; we are not claiming that the
soul is literally governed by Bohmian dynamics. The Noether
analogy is structural; we are not claiming that the soul's life
admits a Lagrangian formulation. These are formal analogies that
illuminate the structure of the claim without committing to a
specific physical realization.

\section{Conclusion}

Chapter 10 has formalized the rizq-conservation claim, established
its structural coherence with modern frameworks, integrated it with
the Akbarian metaphysical substrate, and identified the empirical
program that would be required to move from Level 2 to Level 3.

The chapter has set the template for the rest of Part III. Each
subsequent chapter will follow the same structure: state the claim,
present the objection, formalize using the Part II apparatus,
integrate with the Akbarian framework, draw implications, and
identify what remains open.

The chapter that follows treats the network-theoretic articulation
of the children-bring-rizq claim. The conservation claim of the
present chapter is an individual-scale claim; Chapter 11 will treat
the corresponding household-scale claim and show that it follows
from individual conservation plus specific Islamic structural
conditions.

\begin{summarybox}[Chapter summary]
\textbf{The claim}: for each soul $i$, the lifetime integral
$K_i = \int_0^{T_i} R_i(t)\, dt$ is fixed in advance. The temporal
distribution $R_i(t)$ varies; the integral is invariant.

\textbf{Formalizations}:
\begin{itemize}[leftmargin=1.5em]
    \item Noether-type: invariance under counterfactual life-path
    transformation.
    \item Stochastic-process: bridge structure with deterministic
    integral.
    \item Bohmian: trajectory fully determined by initial wave
    function; motion real.
\end{itemize}

\textbf{Akbarian substrate}: the integral is the soul's structural
capacity for receiving the manifestation of \arb{ar-Razzaq};
\arb{tajalli} provides the dynamic content.

\textbf{Empirical strategy}: natural experiments (twins reared
apart, lottery winners, adverse events, longitudinal cohorts).
None definitively tested; research program is open.

\textbf{Implications}: anxiety reduction is rational; effort
allocation should focus on temporal distribution rather than
integral maximization; charity does not reduce lifetime provision.

\textbf{Level achieved}: Level 2. The structural claim is
formalized; the operational mechanism is not yet specified.

\textbf{What comes next}: Chapter 11 develops the network-theoretic
articulation of the children-bring-rizq claim, building on the
individual conservation result.
\end{summarybox}

\cleardoublepage
